\documentclass[11pt]{article}
\usepackage[T1]{fontenc}
\usepackage[utf8]{inputenc}
\usepackage[english]{babel}
\usepackage{amsmath,amssymb,mathtools,array,booktabs,longtable}
\usepackage[a4paper,margin=1.45cm]{geometry}
\usepackage{microtype}
\setlength{\parindent}{1.25em}
\setlength{\parskip}{0.10em}
\setlength{\emergencystretch}{10em}
\allowdisplaybreaks
\raggedbottom
\sloppy
\hbadness=10000
\vbadness=10000
\hfuzz=2pt
\vfuzz=10pt
\newcommand{\pair}[2]{(#1\,|\,#2)}
\newcommand{\eps}{\varepsilon}
\begin{document}

\section*{\S\ 4. The decomposition identities.}
We now consider the most general case of identities among arbitrary symbol and variable rows in whose final expressions, without performing substitution (11.), a row $p_\tau$ occurs resolved into its individual rows $y^{(1)}\cdots y^{(\tau)}$; we call these identities ``decomposition identities.'' Following the remark at the end of the preceding paragraph, we first determine them for the case of only two symbol rows, that is, we expand the expression $\pair{q_\rho A^\sigma}{B_{\rho+\sigma-\tau}p_\tau}$. We set:
\begin{equation}
\begin{aligned}
q_{\rho-\alpha}^{(i_1\ldots i_\alpha)}&\sim p_{n-\rho}\,y^{(i_1)}\cdots y^{(i_\alpha)},\qquad
p_{\tau-\alpha}^{(i_1\ldots i_\alpha)}=y^{(i_{\alpha+1})}\cdots y^{(i_\tau)},\qquad
p_\tau=y^{(1)}\cdots y^{(\tau)}.
\end{aligned}
\tag{22}
\end{equation}
With respect to the numbers $\rho,\sigma,\tau$ four cases must be distinguished:
\[
\begin{array}{llll}
1)&\rho+\sigma\le n,& \tau\le\rho,& \tau\le\sigma;\\[2pt]
2)&\rho+\sigma\le n,& \tau\ge\rho,& \tau\le\sigma;\\[2pt]
3)&\rho+\sigma\le n,& \tau\le\rho,& \tau>\sigma;\\[2pt]
4)&\rho+\sigma\le n,& \tau\ge\rho,& \tau\ge\sigma.
\end{array}
\]
We shall single out the first case, and then briefly characterize the others from it. By (20.), (10.) and (9.) we obtain, when we decompose the row $y^{(1)}y^{(2)}\cdots y^{(\tau)}B_{\rho+\sigma-\tau}$ into $y^{(1)}$ and $y^{(2)}\cdots y^{(\tau)}B_{\rho+\sigma-\tau}$:
\begin{equation}
\begin{aligned}
\pair{q_\rho A^\sigma}{y^{(1)}y^{(2)}\cdots y^{(\tau)}B_{\rho+\sigma-\tau}}
&=\pair{q_\rho^{(1)}A^\sigma}{y^{(2)}\cdots y^{(\tau)}B_{\rho+\sigma-\tau}}\\
&\quad+(-1)^\rho\pair{q_\rho[A_{n-\sigma}y^{(1)}]^{\sigma-1}}{y^{(2)}\cdots y^{(\tau)}B_{\rho+\sigma-\tau}}.
\end{aligned}
\tag{23}
\end{equation}
The last term in (23.) now has exactly the form of the original expression--only $\sigma$ is replaced by $\sigma-1$, and $\tau$ by $\tau-1$--and therefore again admits equation (23.). Thus, since $\tau\le\sigma$, we can apply operation (23.) $\tau$ times and, taking account of (10.), obtain the relation:
\begin{equation}
\begin{aligned}
\pair{q_\rho A^\sigma}{p_\tau B_{\rho+\sigma-\tau}}
&=(-1)^{\rho\sigma+(\sigma+\tau)(n+1)}
  \pair{q_\rho B^{n-\rho-\sigma+\tau}}{A_{n-\sigma}p_\tau}\\
&\quad+\sum_{i_1=1}^{\tau}(-1)^{\rho(i_1-1)}
 \pair{q_{\rho-1}^{(i_1)}[A_{n-\sigma}y^{(1)}\cdots y^{(i_1-1)}]^{\sigma-i_1+1}}
 {y^{(i_1+1)}\cdots y^{(\tau)}B_{\rho+\sigma-\tau}}.
\end{aligned}
\tag{24}
\end{equation}
Every term under the summation sign is again of the type of the original expression; $\rho$ is replaced by $\rho-1$, $\sigma$ by $\sigma-i_1+1$, and $\tau$ by $\tau-i_1$, so that condition 1) is even more strongly satisfied. Since $\tau\le\rho$, we can therefore apply operation (24.) $\tau$ times and obtain the desired decomposition identity:
\begin{equation}
\begin{aligned}
\pair{q_\rho A^\sigma}{p_\tau B_{\rho+\sigma-\tau}}
&=\eps_{i_0}\pair{q_\rho B^{n-\rho-\sigma+\tau}}{A_{n-\sigma}p_\tau}\\
&\quad+\sum_{i_1=1}^{\tau}\eps_{i_1}
 \pair{q_{\rho-1}^{(i_1)}B^{n-\rho-\sigma+\tau}}{A_{n-\sigma}p_{\tau-1}^{(i_1)}}\\
&\quad+\sum_{i_2=i_1+1}^{\tau}\eps_{i_2}
 \pair{q_{\rho-2}^{(i_1i_2)}B^{n-\rho-\sigma+\tau}}{A_{n-\sigma}p_{\tau-2}^{(i_1i_2)}}+\cdots\\
&\quad+\eps_{i_\tau}\pair{q_{\rho-\tau}^{(i_1\ldots i_\tau)}B^{n-\rho-\sigma+\tau}}{A_{n-\sigma}}\\
&=\sum_{\alpha=0}^{\tau}\sum_{i_\alpha=i_{\alpha-1}+1}^{\tau}\eps_{i_\alpha}
 \pair{q_{\rho-\alpha}^{(i_1\ldots i_\alpha)}B^{n-\rho-\sigma+\tau}}{A_{n-\sigma}p_{\tau-\alpha}^{(i_1\ldots i_\alpha)}}.
\end{aligned}
\tag{25}
\end{equation}
For the sign factor $\eps$ one obtains from (24.):
\begin{equation}
\begin{aligned}
\eps_{i_\alpha}&=(-1)^{\delta_{i_\alpha}},\qquad
\delta_{i_\alpha}=\sum_{\beta=1}^{\alpha}(\rho-\beta+1)(i_{\beta-1}+i_\beta+1)\\
&\quad +(\rho-\alpha)(\sigma+i_\alpha+\alpha)+(\sigma+\tau+\alpha)(n+1),
\end{aligned}
\tag{26}
\end{equation}
where $i_0=0$ is to be put. The summation sign in (25.) is to be understood as follows: to each combination $i_1i_2\ldots i_{\alpha-1}$, where $i_1<i_2<\cdots<i_{\alpha-1}$, all values $i_\alpha$ for which $i_{\alpha-1}<i_\alpha\le\tau$ are adjoined. Taking (26.) into account, one easily verifies that the individual sums in (25.) satisfy the differential equations analogous to (19.):
\begin{equation}
 \left(\frac{\partial}{\partial y^{(i)}}\middle|y^{(i+1)}\right)=0,
 \qquad (i=1,2,\ldots,\tau-1),
\tag{27}
\end{equation}
and hence depend only on the row $p_\tau$. Thus the decomposition identity is an indecomposable identity among the rows $A,B,q,p$; we shall return to its explicit representation in the next paragraph.

In case 2), operation (23.) can likewise be applied $\tau$ times, but operation (24.) stops after the $\rho$-th application. Thus in the final expression of (25.) the first summation sign must be replaced by $\sum_{\alpha=0}^{\rho}$. In cases 3) and 4), operation (23.) can be carried out only $\sigma$ times; in place of (24.) one then has the relation:
\begin{equation}
\begin{aligned}
\pair{q_\rho A^\sigma}{p_\tau B_{\rho+\sigma-\tau}}
&=(-1)^{\rho\sigma}
 \pair{q_\rho^{(\sigma+1\ldots\tau)}B^{n-\rho+\tau-\sigma}}{A_{n-\sigma}p_{\tau-\sigma}^{(\sigma+1\ldots\tau)}}\\
&\quad+\sum_{i_1=1}^{\sigma}(-1)^{\rho(i_1-1)}
 \pair{q_{\rho-1}^{(i_1)}[A_{n-\sigma}y^{(1)}\cdots y^{(i_1-1)}]^{\sigma-i_1+1}}
 {y^{(i_1+1)}\cdots y^{(\tau)}B_{\rho+\sigma-\tau}};
\end{aligned}
\tag{28}
\end{equation}
repeating (28.) $\tau-\sigma$ times, where the summation sign becomes $\sum_{i_\beta=i_{\beta-1}+1}^{\sigma+\beta-1}$--as the $(\sigma-i_1+1)$-fold application of (23.) to the terms under the summation sign in (28.) shows--gives all terms of the individual sum in (25.) corresponding to the index $\alpha=\tau-\sigma$, with the corresponding sign. Then cases 3) and 4) pass over into cases 1) and 2), for the numbers corresponding to $\sigma,\tau$ become $\sigma'=\sigma-i_{\tau-\sigma}+\tau-\sigma$, $\tau'=\tau-i_{\tau-\sigma}=\sigma'$; and as the procedure continues, $\tau'<\sigma'$. Thus in the final expression of (25.) the first summation sign must be replaced by $\sum_{\alpha=\tau-\sigma}^{\tau,\rho}$.

Analogously, from (21.) one obtains the dual decomposition identity, in which a row $q_\rho$ is resolved into the individual rows $v^{(1)},\ldots,v^{(\rho)}$. If we put
\begin{equation}
 \dot p_{\tau-\beta}^{(i_1\ldots i_\beta)}\sim v^{(i_1)}\cdots v^{(i_\beta)}q_{n-\tau};\qquad
 \dot q_{\rho-\beta}^{(i_1\ldots i_\beta)}=v^{(i_{\beta+1})}\cdots v^{(i_\rho)};\qquad
 q_\rho=v^{(1)}\cdots v^{(\rho)},
\tag{29}
\end{equation}
then
\begin{equation}
 \pair{q_\rho A^\sigma}{p_\tau B_{\rho+\sigma-\tau}}
 =\sum_{\beta=\max(0,\tau-\sigma)}^{\min(\tau,\rho)}
   \sum_{i_\beta=i_{\beta-1}+1}^{\rho}\eps
\pair{\dot q_{\rho-\beta}^{(i_1\ldots i_\beta)}B^{n-\rho-\sigma+\tau}}{A_{n-\sigma}\dot p_{\tau-\beta}^{(i_1\ldots i_\beta)}} ,
\tag{30}
\end{equation}
where the summation index $\beta$ runs from the larger of the lower values to the smaller of the upper values. All terms of an individual sum in (25.) and (30.) are polars of a single form, produced by the polar processes
\[
\left(\frac{\partial}{\partial p_{\tau-\alpha}}\middle|p_{\tau-\alpha}^{(i_1\ldots i_\alpha)}\right),\qquad
\left(q_{\rho-\alpha}^{(i_1\ldots i_\alpha)}\middle|\frac{\partial}{\partial q_{\rho-\alpha}}\right).
\]
To express this, we denote by $\Pi$ an aggregate of such polar operations, multiplied by constants, and obtain, in place of (25.) and (30.), the relation
\begin{equation}
 \pair{q_\rho A^\sigma}{p_\tau B_{\rho+\sigma-\tau}}
 =\sum_{\alpha=\max(0,\tau-\sigma)}^{\min(\tau,\rho)}
 \Pi\pair{q_{\rho-\alpha}B^{n-\rho-\sigma+\tau}}{A_{n-\sigma}p_{\tau-\alpha}}.
\tag{31}
\end{equation}

\section*{\S\ 5. The decomposition identities in explicit form.}
If in (25), in case 4), we specialize $\tau$ to $\sigma+\rho$, we obtain
\begin{equation}
 \pair{q_\rho A^\sigma}{y^{(1)}\cdots y^{(\rho+\sigma)}}
 =\sum_{1\le i_1<\cdots<i_\rho\le \rho+\sigma}
 \eps_{i_\rho}\pair{q_\rho}{y^{(i_1)}\cdots y^{(i_\rho)}}
 \pair{A^\sigma}{y^{(i_{\rho+1})}\cdots y^{(i_{\rho+\sigma})}},
\tag{32}
\end{equation}
where $\eps_{i_\rho}=(-1)^{\delta_{i_\rho}}$ and $\delta_{i_\rho}=\rho(\rho+1)/2+i_1+\cdots+i_\rho$. This is a more general form of the product theorem for matrices than (16.) and (17.), and follows by Laplace expansion of the determinant of the product matrix
\[
 \sum\pm(v^{(1)}y^{(1)})\cdots(v^{(\rho)}y^{(\rho)})(a^{(1)}y^{(\rho+1)})\cdots(a^{(\sigma)}y^{(\rho+\sigma)}).
\]
Thus the more general product theorem is composed out of the identities (20.) and (21.), respectively, and the selection of the special form (16.), (17.) is thereby explained.

Relation (32.) now gives the means for explicitly representing the decomposition identities. For this purpose we regard the forms occurring in (31.) as ground forms; that is, we perform substitution (11.), which changes nothing in the explicit representation in the rows $A,B$, since by (8.) the newly introduced rows $R$ can be expressed through the rows $A,B$ themselves and not through their component rows $a^{(1)}a^{(2)}\ldots$, $b^{(1)}b^{(2)}\ldots$. Thus let
\begin{equation}
 \pair{q_{\rho-\alpha}B^{n-\rho-\sigma+\tau}}{A_{n-\sigma}p_{\tau-\alpha}}
 =\pair{R_{\rho-\alpha}p_{n-\rho+\alpha}}{R^{\tau-\alpha}p_{\tau-\alpha}}.
\tag{33}
\end{equation}
Then, for the individual sum of (25.) corresponding to the index $\alpha$, we get
\begin{equation}
\begin{aligned}
&\sum \eps_{i_\alpha}\pair{R_{\rho-\alpha}p_{n-\rho}y^{(i_1)}\cdots y^{(i_\alpha)}}{R^{\tau-\alpha}y^{(i_{\alpha+1})}\cdots y^{(i_\tau)}}\\
&\quad=\sum \eps_{i_\alpha}\pair{[R_{\rho-\alpha}p_{n-\rho}]^\alpha y^{(i_1)}\cdots y^{(i_\alpha)}}{R^{\tau-\alpha}y^{(i_{\alpha+1})}\cdots y^{(i_\tau)}}\\
&\quad=\eps\pair{[R_{\rho-\alpha}p_{n-\rho}]^\alpha R^{\tau-\alpha}}{p_\tau}
 =\eps\pair{R^{\tau-\alpha}q_{n-\tau}}{R_{\rho-\alpha}p_{n-\rho}}.
\end{aligned}
\tag{34}
\end{equation}
Thus an explicit final expression is in fact obtained, and the symmetric form in which $q_\rho$ and $p_\tau$ occur shows that (30.) leads to the same final expression, which is therefore independent of the original decomposition. There is a difference between (25.) and (30.) only as long as the rows $y$ and $v$ have independent meaning, so that the decomposition identity passes, by our definition, into a decomposable identity. To arrive at identities with more symbol rows, one need only, according to \S\ 3 (end), resolve a row $q_\rho$ into $q_{\rho_1}C^{\sigma_1}$, or $p_\tau$ into $p_{\tau_1}C_{\tau-\tau_1}$. The resulting expressions
\begin{equation}
 \pair{q_{\rho_1}C^{\sigma_1}}{[R^{\tau-\alpha}q_{n-\tau}]_\lambda R_{\rho+\sigma_1-\alpha-\lambda}},
 \qquad
 \pair{[R_{\rho-\alpha}p_{n-\rho}]^\lambda R^{\tau-\alpha}}{p_{\tau_1}C_{\tau-\tau_1}}
\tag{35}
\end{equation}
are exactly of the type that occurs for two symbol rows, and therefore again admit explicit representation after a substitution corresponding to (33.). Repetition therefore gives explicit representation for arbitrarily many rows. It may also be noted that the first and last sums in (25.) and (30.), respectively, give the explicit representation without applying (33), solely by formula (32), or else reduce altogether to a single term containing all rows explicitly. The relation obtained from (32.) by resolving $q_\rho$ into $q_{\rho_1}C^{\sigma_1}$ and so forth can be regarded as the product theorem in its most general form.

In summary we shall show:

\emph{Theorem III:} With the identities
\begin{equation}
 \pair{q_\rho A^\sigma}{p_\tau B_{\rho+\sigma-\tau}}
 =\sum_{\alpha=\max(0,\tau-\sigma)}^{\min(\tau,\rho)}
 \eps\pair{R^{\tau-\alpha}q_{n-\tau}}{R_{\rho-\alpha}p_{n-\rho}},
\tag{36}
\end{equation}
where the rows $R$ are determined by (33), together with those obtained from them by resolving $q_\rho,p_\tau$ into $q_{\rho_1}C^{\sigma_1}$, etc., the totality of indecomposable identities is exhausted.

Indeed, the resulting identities are the most general of their kind, since the individual rows are no longer subject to any restriction as to weight or number, as was still the case in (20.) and (21.). Nor do any further identities among these rows exist; for when a row $p_\tau$ is resolved into $y^{(1)}\cdots y^{(\tau)}$, the most general identities are composed out of (20.), hence by \S\ 3 (end) out of (20a.); and the composition discussed following (16.) is precisely the one carried out in \S\ 4. The transition to arbitrarily many rows is supplied, as the discussion of (20a.) shows, by applying always the same operation to the expressions that arise from resolving $q_\rho$ into $q_{\rho_1}C^{\sigma_1}$, and so on. It also follows that the direct transition from (20.) instead of from (20a.) gives nothing new. This is easily verified by calculation, by specializing once a row $q_\rho$ in (25.), and another time carrying out the transition from (20.) by the method of \S\ 4. In both cases the same sums arise, which pass by application of the general product theorem (see above) into the identities obtained from (36.). The identities (20.), (21.) correspond to the special cases $\tau=1$, $\rho=1$; only two individual sums then occur, and hence, in accordance with an earlier remark, explicit representation without the substitution (33.), agreeing with what was found before.

At the end let us briefly mention two special cases of the decomposition identity. Put in (31): $\tau=\sigma$, $\rho=n-\sigma$, and denote the row $p_\tau$ by $p'_\sigma\sim q'_{n-\sigma}$. Then
\begin{equation}
\begin{aligned}
\pair{A^\sigma}{p_\sigma}\pair{B^\sigma}{p'_\sigma}
-\pair{A^\sigma}{p'_\sigma}\pair{B^\sigma}{p_\sigma}
=\sum_{\alpha=1}^{\min(\sigma,n-\sigma)}
\Pi\pair{q_{n-\sigma-\alpha}B^\sigma}{A_{n-\sigma}p_{\sigma-\alpha}}.
\end{aligned}
\tag{37}
\end{equation}
This is Clebsch's formula\footnote{Clebsch, loc. cit., pp. 25--32. Clebsch proves that the individual terms on the right-hand sides depend only on the rows $A,B$; the explicit representation would be obtained by applying (32.) to his expressions $\Phi_h$.} for expanding a form with two cogredient rows $p_\sigma,p'_\sigma$ into elementary covariants. The elementary covariants belonging to a form $\pair{A^\sigma}{p_\sigma}^m\pair{B^\sigma}{p'_\sigma}^n$ are then the following:
\begin{equation}
 \left\{\sum_{\alpha=1}^{\min(\sigma,n-\sigma)}
 \eps\pair{q_{n-\sigma-\alpha}B^\sigma}{A_{n-\sigma}p_{\sigma-\alpha}}
 \right\}^{\rho}
 \pair{A^\sigma}{p_\sigma}^{m-\rho}\pair{B^\sigma}{p'_\sigma}^{n-\rho},
 \qquad (\rho=0,1,\ldots,m,n).
\tag{38}
\end{equation}
They arise, as we shall briefly say, from the ground form by ``cogredient folding of defect $\rho$'' (cf. \S\ 2).

Second, put $\tau=\sigma-\lambda$, $\rho=n-\sigma-\lambda$, $B^\sigma=A^\sigma$. Then
\begin{equation}
\begin{aligned}
\pair{q_{n-\sigma-\lambda}A^\sigma}{A_{n-\sigma}p_{\sigma-\lambda}}
&=(-1)^\lambda\pair{q_{n-\sigma-\lambda}A^\sigma}{A_{n-\sigma}p_{\sigma-\lambda}}\\
&\quad+(-1)^{(n-\sigma)(\sigma-\lambda)}
 \sum_{\alpha=1}^{\min(\sigma-\lambda,n-\sigma-\lambda)}
 \Pi\pair{q_{n-\sigma-\lambda-\alpha}A^\sigma}{A_{n-\sigma}p_{\sigma-\lambda}}.
\end{aligned}
\tag{39}
\end{equation}
Thus, as noted in \S\ 2, the conditions (12.) characterizing the symbol rows for odd defect $\lambda$ are consequences of the conditions of higher defect. For a linear form $\pair{A^\sigma}{p_\sigma}=\pair{B^\sigma}{p_\sigma}$, (39.) implies that its irreducible invariant formations that are quadratic in the coefficients reduce to ones of even defect; this agrees with the known fact that a linear form in the binary and ternary domains has no invariant formations.

It should be noted that the general identities were originally derived under the assumption that the individual rows satisfy conditions (12.). Since, however, these individual rows enter the identities only linearly, the latter also hold for the more general rows that do not satisfy conditions (12.).

\section*{\S\ 6. Explicit representation of the invariant formations.\\Folding and differentiation processes.}
The results of the preceding paragraphs allow us to complete the theorem stated in \S\ 2 on symbolic representability (the first fundamental theorem) by adding explicit representability; namely, we show:

\emph{Theorem IV:} ``All invariant formations of arbitrary ground forms can be represented explicitly by matrix products; that is, aside from variable rows $p_\rho$, only the rows $S^\sigma,\ldots,T^\tau$ of the original forms, or the rows $P,Q,R$ derived from them by substitution (9.) or (11.), enter the final expressions.''

For the proof, suppose that an invariant formation depends, besides variable rows, on the rows $S^\sigma,\ldots,T^\tau$; and suppose that these rows have been resolved sufficiently far into individual rows $S^\sigma=S^{\sigma_1}\cdots S^{\sigma_i},\ldots,T^\tau=T^{\tau_1}\cdots T^{\tau_k}$ to make representation by determinants possible. Let the rows be ordered in a sequence $S^\sigma,\ldots,T^\tau$ that is arbitrary in itself but fixed once and for all. We pick out an aggregate of determinants that depends on the first total row $S^\sigma=S^{\sigma_1}\cdots S^{\sigma_i}$, not merely on individual component rows. This dependence, however, is by \S\ 3--\S\ 5 a consequence of the general identities. If we now resolve a variable row $p_\rho$ into the rows $y,v$, or one of the later rows $T$ into the individual rows $t$, respectively $\tau$, then the most general identities are composed out of (20.) and (21.). But these latter identities always lead to an expression containing the row $S^\sigma=S^{\sigma_1}\cdots S^{\sigma_i}$ explicitly, namely the left-hand side of (20.) and (21.); for by (19.) one verifies that only the complete sum on the right--not a part of it corresponding to a term of (20a.)--has the property of depending only on $S^\sigma$. As the procedure continues, all rows that have once entered explicitly remain explicit; application of the identities can at most require the combining of several rows into a single one, i.e. substitution (9.). If finally only a single row $T^\tau$ remains to be brought to explicit form, two cases are possible. There may still occur a free variable row, that is, one not connected with other rows by substitution (9.); resolving it into component rows $y$ or $v$ completes the procedure and gives, as in (31.), polars of one or more forms containing all rows explicitly. If no free variable row remains, then from the manner in which they arose we recognize, in the totality of the expressions that contain the individual rows $t$ or $\tau$ but depend on the row $T^\tau$, the terms of a decomposition identity; by \S\ 5, however, these always admit explicit representation in all rows after applying substitution (11.). This proves the theorem in full generality. It should also be noted that when only two symbol rows occur, substitution (11.) never appears. For, since $\sigma,n-\sigma,\tau,n-\tau<n$, no variable rows can enter only when the two symbol rows are united in a single determinant and thus occur explicitly; but as soon as variable rows enter, (34.) and (33.) show that substitution (11.) becomes unnecessary.

It follows from what has just been said that, in order to generate all invariant formations, it suffices to unite the symbol rows, after adjoining variable rows, into all possible matrix products. The most general matrix product is now of the form
\begin{equation}
 \pair{P^{\mu_1}\cdots P^{\mu_i}}{Q_{\nu_1}\cdots Q_{\nu_k}},\qquad \sum\mu=\sum\nu,
\tag{40}
\end{equation}
where the $P$ and $Q$ may be rows of the original forms, or may have arisen from the original rows by a sequence of substitutions (9.) or (11.), or finally may be variable rows. But the expressions (40.) can be generated by successively uniting two symbol rows at a time into a matrix product, with the help of variable rows; for from
\[
 \pair{q_{n-\mu-\lambda}P^\mu}{Q_{\nu_1}p_{n-\mu-\nu_1-\lambda}}=\pair{R_\lambda Q_{\nu_1}}{p_{n-\mu-\nu_1-\lambda}}
\]
one obtains, by uniting with the row $Q_{\nu_2}$, an expression
\[
 \pair{[R_\lambda Q_{\nu_1}]^{n-\lambda-\nu_1}}{Q_{\nu_2}p_{n-\mu-\nu_1-\nu_2-\lambda}}
 =\pair{q_{n-\mu-\lambda}P^\mu}{Q_{\nu_1}Q_{\nu_2}p_{n-\mu-\nu_1-\nu_2-\lambda}},
\]
and hence, by continuing the procedure, an expression (40.). If $P$ and $Q$ are not rows of the original forms, then (9.) and (11.), in conjunction with what was just shown, imply that they too arose through a sequence of unions of two symbol rows at a time. Thus:

\emph{Theorem V:} The process for generating all invariant formations, the folding process, consists in successively uniting two symbol rows at a time into a matrix product, with the help of variable rows.

From two symbol rows $S^\sigma,T^\tau$, where $\sigma\ge\tau$, the following matrix products can be formed:
\begin{align}
&\pair{q_{n-\sigma-\lambda}S^\sigma}{T_{n-\tau}p_{\tau-\lambda}}, &&(\lambda=1,2,\ldots,\min(\tau,n-\sigma)),\tag{41}\\
&\pair{q_{n-\tau-\mu}T^\tau}{S_{n-\sigma}p_{\sigma-\mu}}, &&(\mu=\sigma-\tau+\lambda),\tag{42}\\
&\pair{q_{n-\tau-\nu}T^\tau}{S_{n-\sigma}p_{\sigma-\nu}}, &&(\nu=1,2,\ldots,\sigma-\tau).
\tag{43}
\end{align}
From (31.), however, the following values are obtained for (42.) and (43.):
\begin{align}
&\pair{q_{n-\sigma-\lambda}S^\sigma}{T_{n-\tau}p_{\tau-\lambda}}
 +\sum_{\alpha}\Pi\pair{q_{n-\sigma-\lambda-\alpha}S^\sigma}{T_{n-\tau}p_{\tau-\lambda-\alpha}},\tag{42a}\\
&\Pi(q_{n-\sigma}S^\sigma)(T_{n-\tau}p_\tau)
 +\sum_{\beta}\Pi\pair{q_{n-\sigma-\beta}S^\sigma}{T_{n-\tau}p_{\tau-\beta}}.
\tag{43a}
\end{align}
Thus the forms (42.) can be expressed linearly by (41.) and polars of (41.), and the forms (43.) by polars of the ground form and of the forms (41.). Likewise (31.) gives the linear dependence of the forms (41.) on (42.) and polars of (42.). Hence, according to Clebsch's definition,\footnote{Clebsch, loc. cit., \S\ 7.} the forms (42.) are equivalent to the forms (41.), while (43.) leads back to the ground form. The generating process is therefore given with equal right by (41.) or (42.); we shall define the process (41.), which is simpler with respect to the number $\lambda$, as the folding process. The number $\lambda$, the defect of the folding (cf. \S\ 2), gives the number of rows missing from the determinant product, namely in that matrix product which, for the given folding, contains the maximal number of rows. Since $\lambda$ runs only up to the smaller of the two values $\tau,n-\sigma$, where $\sigma\ge\tau$, we have
\begin{equation}
 \lambda\le\frac n2,\ n\text{ even};\qquad
 \lambda\le\frac{n-1}{2},\ n\text{ odd}.
\tag{44}
\end{equation}
The reduction of (42.) and (43.) to (41.) remains valid even when, through further folding, the rows $p,q$ have passed into symbol rows; for by (36.) one has
\[
 \pair{A^{n-\tau-\varkappa}T^\tau}{S_{n-\sigma}B_{\sigma-\varkappa}}
 =\sum_{\alpha=\max(0,\varkappa-\sigma+\tau)}^{\min(\tau,n-\sigma)}
 \eps\pair{R^{\tau-\alpha}B^{n-\sigma+\varkappa}}{A_{\tau+\varkappa}R_{n-\sigma-\alpha}},
\]
where the rows $R$ have arisen from $S$ and $T$ according to (33.). Thus these forms containing symbol rows only are obtained by folding $\pair{q_\tau A^{n-\tau-\varkappa}}{B_{\sigma-\varkappa}p_{n-\sigma}}$ or $\pair{q_{\tau-\alpha}B^{n-\sigma+\varkappa}}{A_{\tau+\varkappa}p_{n-\sigma-\alpha}}$--according as $\sigma-\tau\gtreqless 2\varkappa$--with the ground form and the forms (41.).

From the symbolic folding process one obtains, in the usual way, the invariant differentiation processes if one merely replaces the symbol rows $S^\sigma,T_{n-\tau}$ by the rows of differentiation symbols $\frac{\partial}{\partial p_\sigma}$, $\frac{\partial}{\partial q_{n-\tau}}$; as is well known, the product $\frac{\partial}{\partial p_{i_1\ldots i_\sigma}}\frac{\partial}{\partial q_{j_1\ldots j_{n-\tau}}}$ then has the real meaning $\frac{\partial^2}{\partial p_{i_1\ldots i_\sigma}\partial q_{j_1\ldots j_{n-\tau}}}$. Since the rows $\frac{\partial}{\partial p_\sigma}$ and $\frac{\partial}{\partial q_{n-\tau}}$ are cogredient with the rows $S^\sigma,T_{n-\tau}$, they satisfy the same identities as these, in particular also the ones existing between (42.) and (42a.), and between (43.) and (43a.). In summary we obtain:

\emph{Theorem VI:} All invariant formations of arbitrary ground forms arise from these by repeated application of the symbolic folding process (41.), or also of the invariant differentiation processes
\begin{equation}
 \left(q_{n-\sigma-\lambda}\frac{\partial}{\partial p_\sigma}\middle|\frac{\partial}{\partial q_{n-\tau}}p_{\tau-\lambda}\right),
 \qquad(\sigma\ge\tau,\ \lambda=1,2,\ldots,\min(\tau,n-\sigma)).
\tag{45}
\end{equation}
It remains to note that in every folding (41.) the total weight of the form, that is, the sum of the weights of the individual variable rows (cf. \S\ 1), decreases by the positive number $2(\lambda^2+\lambda(\sigma-\tau))$. This implies, independently of the general finiteness proof, the finiteness of the form row, that is, of the totality of invariant formations of a given ground form that are linear in the coefficients.\footnote{Cf. Clebsch, loc. cit., \S\ 11 and 12: finiteness of the system of elementary covariants.}

It should further be noted that Theorem VI still holds for forms that do not satisfy conditions (12.). For each factor in (3), $\pair{S^\rho}{p_\rho}^{m_\rho}$, can be replaced by a product of $m_\rho$ linear factors $\pair{A^\rho}{p_\rho}\pair{B^\rho}{p_\rho}\cdots\pair{M^\rho}{p_\rho}$; the invariant formations thereby pass into those of a system of linear forms, which need only afterwards be set equal to one another. For every single linear form, however, conditions (12.)--which, when differential symbols are introduced, contain only second-order differential quotients--are always satisfied.

\section*{\S\ 7. Expansion of the general forms according to normal forms.}
In what follows (\S\ 8) we shall derive a principal property of the folding process (41.), namely its composability from fundamental foldings. For this we need the transfer to the $n$-ary domain of an important theorem given by Mertens\footnote{F. Mertens, \emph{\"Uber invariante Gebilde quatern\"arer Formen} (see the Introduction), cited as ``Mertens.''} in the quaternary domain. This theorem says that any finite system of forms can be replaced by an equivalent system of normal forms. By a ``normal form'' we mean--generalizing the binary and ternary definition\footnote{Study, loc. cit., p. 54.}--a form all of whose invariant formations linear in the coefficients vanish identically, with the exception of the ground form itself. Thus by Theorem VI (\S\ 6) a normal form $\varphi$ satisfies the system of differential equations
\begin{equation}
 \left(q_{n-\sigma-\lambda}\frac{\partial}{\partial p_\sigma}\middle|\frac{\partial}{\partial q_{n-\tau}}p_{\tau-\lambda}\right)\varphi=0,
 \qquad(\sigma\ge\tau,\ \lambda=1,2,\ldots,\min(\tau,n-\sigma)),
\tag{46}
\end{equation}
where the rows $q_{n-\sigma-\lambda},p_{\tau-\lambda}$ are to be regarded as arbitrary quantities. In the quaternary domain--taking (39.) into account for $\sigma=\tau=2$--these differential equations agree completely with the ten differential equations given by Mertens without introduction of the arbitrary rows $p,q$ (loc. cit., formula 28). Their knowledge, together with Theorem VI, allows us to transfer his proof almost word for word to $n$ variables. The only thing to add is the verification, obtained by the decomposition identity (30.), that the constructed forms are normal forms; in the quaternary domain this verification still follows without further transformation. It is therefore enough to give a short exposition of the proof, and for simplicity in the case relevant below, namely that of a single ground form $f$ and its form row (cf. \S\ 6, end); for invariant formations of arbitrary degree in the coefficients of arbitrarily many ground forms, the proof remains the same.

The basic idea of the proof is to replace, by introducing the transformed coefficients, a form with $n-1$ rows $p_\rho$ by forms with $n$ rows $\xi$ cogredient to the $x$, namely the rows of the transformation quantities. From these forms the sought normal forms are obtained directly by invariant differentiation processes. The expansion of the general forms according to normal forms is mediated by a series expansion for forms with $\rho$ ($\rho<n$) cogredient rows $\xi$; invariant differentiation processes lead back to the forms in the original rows $p_\rho$.

Let $\xi^{(1)},\ldots,\xi^{(n)}$ be the rows of transformation quantities, so that
\begin{equation}
\begin{aligned}
x_i&=\xi_i^{(1)}x'_1+\xi_i^{(2)}x'_2+\cdots+\xi_i^{(n)}x'_n,\\
p_{i_1\ldots i_\rho}&=\sum_j(\xi_{i_1}^{(j_1)}\cdots \xi_{i_\rho}^{(j_\rho)})p'_{j_1\ldots j_\rho}.
\end{aligned}
\tag{47}
\end{equation}
The transformed coefficients of the forms $f,\varphi,\ldots$ will be denoted by $(f),(\varphi),\ldots$, and let
\[
 f=\pair{S^1}{p_1}^{m_1}\pair{S^2}{p_2}^{m_2}\cdots\pair{S^{n-1}}{p_{n-1}}^{m_{n-1}}.
\]
Then
\begin{equation}
\begin{aligned}
(f)&=\pair{S^1}{\xi^{(1)}}^{m_{11}}\cdots\pair{S^1}{\xi^{(n)}}^{m_{1n}}
\pair{S^2}{\xi^{(1)}\xi^{(2)}}^{m_{21}}\cdots
\pair{S^{n-1}}{\xi^{(2)}\cdots\xi^{(n)}}^{m_{n-1,n}}\\
&\quad\cdot\pair{s^{(1)}}{\xi^{(1)}}^{k_1}\pair{s^{(2)}}{\xi^{(2)}}^{k_2}\cdots\pair{s^{(n)}}{\xi^{(n)}}^{k_n},
\end{aligned}
\tag{48}
\end{equation}
where $m_{11}+\cdots+m_{1n}=m_1$, and so on. From each coefficient $(f)$ for which
\begin{equation}
 k_1\ge k_2\ge k_3\ge\cdots\ge k_n,
\tag{49}
\end{equation}
the differentiation process exhausting exactly the rows $\xi$,
\begin{equation}
\left(\frac{\partial}{\partial\xi^{(1)}}\middle|p_1\right)^{k_1-k_2}
\left(\frac{\partial}{\partial\xi^{(1)}}\frac{\partial}{\partial\xi^{(2)}}\middle|p_2\right)^{k_2-k_3}
\cdots
\left(\frac{\partial}{\partial\xi^{(1)}}\cdots\frac{\partial}{\partial\xi^{(n-1)}}\middle|p_{n-1}\right)^{k_{n-1}-k_n}
\left(\frac{\partial}{\partial\xi^{(1)}}\cdots\frac{\partial}{\partial\xi^{(n)}}\right)^{k_n}
\tag{50}
\end{equation}
gives the form
\begin{equation}
 \varphi=\bigl(s^{(1)}\mid p_1\bigr)^{k_1-k_2}
 \bigl(s^{(1)}s^{(2)}\mid p_2\bigr)^{k_2-k_3}\cdots
 \bigl(s^{(1)}\cdots s^{(n-1)}\mid p_{n-1}\bigr)^{k_{n-1}-k_n}
 \bigl(s^{(1)}s^{(2)}\cdots s^{(n)}\bigr)^{k_n}.
\tag{51}
\end{equation}
The forms $\varphi$ are the sought normal forms. Indeed, they have the following properties defining normal forms.

1. They are invariant formations of the ground form $f$ that are linear in the coefficients. This follows from the fact that they arise from $f$ by the invariant differentiation processes $\left(\frac{\partial}{\partial p_\rho}\middle|\xi^{(i_1)}\xi^{(i_2)}\cdots\xi^{(i_\rho)}\right)$ and (50.). Thus (by \S\ 6, end) they can be expressed linearly by polars of the finite form row of $f$. These polars are obtained by regarding the variables $p_\rho$ entering $\varphi$ as coefficients of a form decomposed into linear forms with the variable rows $p_\rho$:
\begin{equation}
 H=(q_{n-1}\mid p_{n-1})^{k_1-k_2}(q_{n-2}\mid p_{n-2})^{k_2-k_3}\cdots(q_1\mid p_1)^{k_{n-1}-k_n}.
\tag{52}
\end{equation}
The simultaneous invariants of the form rows $f$ and $H$ give all relevant polars, since the latter must be of the same dimension in the individual rows $p_\rho$ as $\varphi$ itself.

2. They satisfy the differential equations (46.). If the differential process (45.) is applied to $\varphi$, then
\[
 \varphi'=\bigl(q_{n-\sigma-\lambda}s^{(1)}\cdots s^{(\sigma)}\mid [s^{(1)}\cdots s^{(\tau)}]_{n-\tau}p_{\tau-\lambda}\bigr),\qquad (\sigma>\tau).
\]
From the decomposition identity (30.) it follows, if the rows $q_\sigma=s^{(1)}\cdots s^{(\sigma)}$, $p_{n-\tau}=[s^{(1)}\cdots s^{(\tau)}]_{n-\tau}$ formed out of the rows $s$ are identified with the $q_\rho,p_\tau$ occurring there, that
\[
 \varphi'=\sum_{\beta=\sigma-\tau+\lambda}^{n-\tau,\sigma}\sum_{i_\beta}\eps\,
 \pair{q_{\sigma-\beta}^{(i_1\ldots i_\beta)}q_{n-\tau+\lambda}}{p_{\sigma+\lambda}p_{n-\tau-\beta}^{(i_1\ldots i_\beta)}}.
\]
Here
\[
 p_{n-\tau-\beta}^{(i_1\ldots i_\beta)}\sim s^{(1)}\cdots s^{(\tau)}s^{(i_1)}\cdots s^{(i_\beta)},
\]
where $i_1,
\ldots,i_\beta$ denote any $\beta$ distinct numbers chosen from the numbers 1 to $\sigma$. Since $\beta>\sigma-\tau$, it follows that among these numbers $i_1,
\ldots,i_\beta$ there must be some between 1 and $\tau$; hence $p_{n-\tau-\beta}^{(i_1\ldots i_\beta)}$ vanishes identically, and therefore every individual matrix product, and with it $\varphi'$, vanishes.

For the equivalence of the system of normal forms with the form row $f$, it remains only to show that all forms of the form row can be expanded in polars of the normal forms in the sense fixed under 1. Let $\Pi$ be an aggregate of polar operations
\begin{equation}
 \left(\frac{\partial}{\partial\xi^{(i)}}\middle|\xi^{(k)}\right),
 \qquad (i\ne k;\ i=1,2,\ldots,\rho;\ k=1,2,\ldots,\rho),
\tag{53}
\end{equation}
and let $\Pi^\sigma$ be an aggregate of the special operations (53.) for which $i=\sigma$ and $k<\sigma$. For a form $F$ of the $\rho$ rows $\xi^{(1)},\ldots,\xi^{(\rho)}$ that has degree $k_\rho$ in the row $\xi^{(\rho)}$, the expansion\footnote{F. Mertens, \emph{\"Uber eine Formel der Determinantentheorie}. Sitzungsberichte der Akademie der Wissenschaften, Vienna, Math.-Naturw. Kl., vol. 91, 2nd section (1885), formula 20).} holds:
\begin{equation}
 F=\sum_{\alpha=0}^{k_\rho}\Pi F_\alpha,
 \qquad(\rho\le n),
\tag{54}
\end{equation}
where $F_\alpha$ has degree $\alpha$ in the last row $\xi^{(\rho)}$ and arises from $F$ by the invariant differential operations
\begin{equation}
 \left(\frac{\partial}{\partial\xi^{(1)}}\cdots\frac{\partial}{\partial\xi^{(\rho)}}\middle|\xi^{(1)}\cdots\xi^{(\rho)}\right)^\alpha
\tag{55}
\end{equation}
and $\Pi^\rho$. If, in constructing $F_\alpha$, one replaces process (55.) by
\begin{equation}
 \left(\frac{\partial}{\partial\xi^{(1)}}\cdots\frac{\partial}{\partial\xi^{(\rho)}}\middle|p_\rho\right)^\alpha,
\tag{56}
\end{equation}
then a form $F_\alpha(p_\rho)$ with only $\rho-1$ rows $\xi^{(1)},\ldots,\xi^{(\rho-1)}$ is obtained, to which (54.) can again be applied. Continuing the procedure leads to forms $G(p_\rho,p_{\rho-1},\ldots,p_1)$. In order to obtain from these the expansion of $F$ according to (54.), one replaces the individual rows $p_\rho$ by $\xi^{(1)},\ldots,\xi^{(\rho)}$ and applies to the resulting forms the polar process $\Pi$ (53.); this gives (cf. (48.)) a sum of transformed coefficients $(G)$. For $\rho=n$, the process (55.) remains at the first step. If $c$ denotes numerical constants and we put, for brevity,
\begin{equation}
 (\xi^{(1)}\xi^{(2)}\cdots\xi^{(n)})\sim\Delta;
 \qquad
 \left(\frac{\partial}{\partial\xi^{(1)}}\frac{\partial}{\partial\xi^{(2)}}\cdots\frac{\partial}{\partial\xi^{(n)}}\right)=\nabla,
\tag{57}
\end{equation}
then, in the case $\rho=n$,
\begin{equation}
 F=\sum c\Delta^\alpha(G),
\tag{58}
\end{equation}
whereas for $\rho<n$ only $\alpha=0$, and hence $\sum c(G)$, appears on the right.

Applying (58.) to a transformed coefficient $(f)$ (48.), and using the commutativity of the polar processes $\Pi^\sigma$ with all operations (56.) for which $\rho\ge\sigma$, as well as the fact that polars of transformed coefficients again give transformed coefficients, one obtains
\[
 G=\left(\frac{\partial}{\partial\xi^{(1)}}\middle|p_1\right)^{\beta_1}
 \left(\frac{\partial}{\partial\xi^{(1)}}\frac{\partial}{\partial\xi^{(2)}}\middle|p_2\right)^{\beta_2}
 \cdots
 \left(\frac{\partial}{\partial\xi^{(1)}}\cdots\frac{\partial}{\partial\xi^{(n-1)}}\middle|p_{n-1}\right)^{\beta_{n-1}}
 \nabla^{\beta_n}\Pi(f)=\sum c\varphi,
\]
and from (58.) it follows that
\begin{equation}
 (f)=\sum c\Delta^\alpha(\varphi)\qquad\text{(Mertens, formula 23)).}
\tag{59}
\end{equation}
To pass from (59.) to the expansion of $f$ itself, we again regard the variables $p_\rho$ as coefficients of a form $H$ (52.) with the degrees $m_1,
\ldots,m_{n-1}$ corresponding to $f$, and put $m=m_1+\cdots+m_{n-1}$. Then, by the definition of invariants,
\[
 f\cdot\Delta^m=\sum(f)(H)=\sum c\Delta^\alpha(\varphi)(H),
\]
and application of the process $\nabla^m$ that exactly exhausts the rows $\xi$ gives
\begin{equation}
 f=\sum c\Phi,
\tag{60}
\end{equation}
where the forms $\Phi$ are derived from $\varphi$ and $H$ by invariant differentiation processes with respect to the variables, and hence are simultaneous invariants of the form rows $\varphi$ and $H$.\footnote{Instead of the direct conclusion, Mertens, \S\ 7, introduces a series of further differential processes that essentially serve to form the form row $H$; this is done by our Theorem VI.} But since $\varphi$ is a normal form, the form row $\varphi$ reduces to the form $\varphi$ itself; the form row $H$ contains all possible invariant variable combinations. Thus the simultaneous invariants of $\varphi$ and $H$ are the polars of $\varphi$ in the sense indicated under 1.; in other words, (60.) gives the expansion of $f$ in polars of the normal forms. It remains to verify the same expansion for an arbitrary form $\theta$ of the form row $f$. Let $\Gamma$ be the normal forms derived from $(\theta)$ by (50.), so that
\begin{equation}
 (\theta)=\sum c\Delta^\beta(\Gamma).
\tag{61}
\end{equation}
The forms $\theta$ are characterized by the relation, valid for every transformed coefficient,
\[
 (\theta)\Delta^\sigma=\sum c(f)\qquad\text{(Mertens, formula 9)).}
\]
But every form $F$ of the $n$ rows $\xi$ which contains the factor $\Delta^\sigma$ also contains the second $\nabla^\sigma\Pi F$ (Mertens, formula 20)); hence
\[
 (\theta)=\sum c\nabla^\sigma(f),\qquad\text{therefore}\qquad \Gamma=\sum c\varphi,
\]
and from (61.) we obtain the relation corresponding to (59.),
\[
 (\theta)=\sum c\Delta^\beta(\varphi),
\]
which implies for $\theta$ the expansion (60.) in polars of the normal forms $\varphi$. Thus:

\emph{Theorem VII.} Every form row $f$ can be replaced by an equivalent system of normal forms.

\section*{\S\ 8. Reduction of the folding process to fundamental foldings.}
Following Theorem VII, we now carry out the reduction of the folding process to fundamental foldings mentioned at the beginning of \S\ 7. We call (cf. \S\ 5) every folding of two cogredient rows $S^\rho,T^\rho$ a ``cogredient folding'', and specifically the cogredient folding of defect 1 of two rows $S^\rho,T^\rho$ a ``fundamental folding of type $\rho$''. We then prove, completing Theorem VI:

\emph{Theorem VIII:} The general folding process (41.) can be composed from the $(n-1)$ fundamental foldings of type $\rho$, where $\rho=1,2,\ldots,n-1$. In other words, to generate all invariant formations it suffices to apply successively the $n-1$ fundamental foldings.

In the binary domain the folding process (41.) coincides directly with the only possible fundamental folding; in the ternary domain I proved Theorem VIII in \S\ 1 of my dissertation (cf. the Introduction). There, because of (44.), the general cogredient foldings are still identical with the fundamental foldings, and it is noteworthy that Theorem VIII for $n$ variables gives the reduction to fundamental foldings, not merely the much weaker reduction to cogredient foldings. The proof is modeled in principle on the ternary proof: the most general foldings are constructed from the fundamental foldings, using the symbolic identities. We shall generate:
\begin{enumerate}
\item arbitrary foldings of defect 1 by cogredient foldings of defect 1, that is, by fundamental foldings (the analogue of the ternary case),
\item arbitrary foldings of defect $\lambda$ by cogredient foldings of defect $\lambda$ and arbitrary foldings of defect 1,
\item cogredient foldings of defect $\lambda$ by cogredient foldings of defect $\lambda-1$ and arbitrary foldings of defect 1.
\end{enumerate}
As is essential, we assume by Theorem VII that the two forms $S$ and $T$ to be folded are normal forms. Thus, by (46.), the relations hold:
\begin{equation}
\begin{aligned}
\pair{q_{n-\sigma_1-\lambda}S^{\sigma_1}}{S_{n-\sigma_2}p_{\sigma_2-\lambda}}&=0,
&& (\sigma_1\ge\sigma_2,\ \lambda=1,2,\ldots,\sigma_2,n-\sigma_1),\\
\pair{q_{n-\tau_1-\lambda}T^{\tau_1}}{T_{n-\tau_2}p_{\tau_2-\lambda}}&=0,
&& (\tau_1\ge\tau_2,\ \lambda=1,2,\ldots,\tau_2,n-\tau_1).
\end{aligned}
\tag{62}
\end{equation}
It is further sufficient, by \S\ 2 (end), to suppose that the forms $S$ and $T$ are linear in all variable rows; that is, the constructed forms belong to the form row
\[
 f=\pair{S^1}{p_1}\pair{S^2}{p_2}\cdots\pair{S^{n-1}}{p_{n-1}}
   \pair{T^1}{p_1}\pair{T^2}{p_2}\cdots\pair{T^{n-1}}{p_{n-1}}.
\]
1) Generation of $\pair{q_{n-\sigma-1}S^\sigma}{T_{n-\tau}p_{\tau-1}}$, where $\sigma>\tau$, from fundamental foldings of type $\tau+\alpha$, $\alpha=0,1,\ldots,\sigma-\tau$. From the fundamental folding of type $\tau$
\[
 \pair{q_{n-\tau-1}S^\tau}{T_{n-\tau}p_{\tau-1}}
\]
we obtain, by the substitution $w\sim T_{n-\tau}p_{\tau-1}$, a form of the form row $f$ with three rows of upper weight index (cf. \S\ 1) $\tau+1$:
\begin{equation}
 \pair{wS^\tau}{p_{\tau+1}}\pair{S^{\tau+1}}{p_{\tau+1}}\pair{T^{\tau+1}}{p_{\tau+1}}.
\tag{63}
\end{equation}
A fundamental folding of type $\tau+1$ applied to (63.) gives
\[
 \pair{q_{n-\tau-2}wS^\tau}{S_{n-\tau-1}p_\tau};
\]
from (31.), by the substitution $q_{n-\tau-2}w=q'_{n-\tau-1}$, this has the value
\[
 \eps\pair{q_{n-\tau-2}wS^{\tau+1}}{S^\tau p_\tau}
 +\sum_{\alpha=1}^{\tau,n-\tau-1}\Pi\pair{q'_{n-\tau-1-\alpha}S^{\tau+1}}{S_{n-\tau}p_{\tau-\alpha}}.
\]
The expression under the summation sign vanishes identically by (62.); thus we have reached a form
\[
 \pair{wS^{\tau+1}}{p_{\tau+2}}\pair{S^{\tau+2}}{p_{\tau+2}}\pair{T^{\tau+2}}{p_{\tau+2}},
\]
which is obtained from (63.) by changing $\tau$ into $\tau+1$. Repeating the process $\sigma-\tau$ times therefore leads to the desired folding:
\[
 \pair{wS^\sigma}{p_{\sigma+1}}=\eps\pair{q_{n-\sigma-1}S^\sigma}{T_{n-\tau}p_{\tau-1}}.
\]
We indicate the process just carried out by the formula
\begin{equation}
 T^\tau S^\tau,
\quad S^{\tau+1},\quad S^{\tau+2},\ldots,S^\sigma;
\tag{64}
\end{equation}
and see that only the normal-form property of $S$, not that of $T$, is used. A process dual to (64.) can be performed, using only the normal-form property of $T$, in the reverse order, namely:
\begin{equation}
 S^\sigma T^\sigma,
\quad T^{\sigma-1},\quad T^{\sigma-2},\ldots,T^\tau,
\tag{65}
\end{equation}
according to the relations:
\begin{align*}
\pair{q_{n-\sigma-1}S^\sigma}{T_{1-\sigma}p_{\sigma-1}}&=\eps\pair{T_{n-\sigma}y}{p_{\sigma-1}},\quad y\sim q_{n-\sigma-1}S^\sigma,\\
\pair{q_{n-\sigma}T^{\sigma-1}}{T_{n-\sigma}yp_{\sigma-2}}&=\eps\pair{T_{n-\sigma+1}y}{p_{\sigma-2}},\\
&\vdots\\
\pair{q_{n-\tau-1}T^\tau}{T_{n-\tau-1}yp_{\tau-1}}&=\eps\pair{q_{n-\sigma-1}S^\sigma}{T_{n-\tau}p_{\tau-1}}.
\end{align*}
We should also point out the connection with the decrease of total weight given in \S\ 6 (end): for foldings of defect 1 this decrease has the value $2(1+(\sigma-\tau))$, and for fundamental foldings the value $2\cdot1$. Therefore, if composition out of fundamental foldings is possible at all, exactly $\sigma-\tau+1$ such foldings are necessarily required, as was carried out above.

2) Generation of general foldings from cogredient and fundamental foldings. The weight decrease for a general folding (41) is $2(\lambda^2+\lambda(\sigma-\tau))=2[\lambda^2+(\sigma-\tau)(1+(\lambda-1))]$. This gives the possibility of decomposing it into a cogredient folding of defect $\lambda$ (weight decrease $2\lambda^2$) and $\sigma-\tau$ foldings of defect 1 with difference of weight indices $\lambda-1$ (weight decrease $2\{1+(\lambda-1)\}$). In the case $\lambda=1$ this decomposition agrees with the one given under 1); we show that it can indeed be realized.

From the cogredient folding of defect $\lambda$
\[
 \pair{q_{n-\tau-\lambda}S^\tau}{T_{n-\tau}p_{\tau-\lambda}}
\]
we obtain, by the substitution $R^\lambda\sim T_{n-\tau}p_{\tau-\lambda}$, a form with the two rows of upper weight indices $\tau+\lambda$ and $\tau+1$:
\begin{equation}
 \pair{R^\lambda S^\tau}{p_{\tau+\lambda}}\pair{S^{\tau+1}}{p_{\tau+1}}.
\tag{66}
\end{equation}
The row with the lower weight index $\tau+1$ belongs to a normal form; hence (66.) admits a folding of defect 1 composed out of $\lambda$ fundamental foldings, in the order (65):
\[
 (R^\lambda S^\tau)S^{\tau+\lambda},\quad S^{\tau+\lambda-1},\quad S^{\tau+\lambda-2},\ldots,S^{\tau+1}.
\]
Taking account of (31.) and (62.), this gives
\[
 \pair{q_{n-\tau-\lambda-1}R^\lambda S^\tau}{S_{n-\tau-1}p_\tau}=\eps\pair{R^\lambda S^{\tau+1}}{p_{\tau+\lambda+1}},
\]
that is, a form obtained from (66.) by replacing $\tau$ by $\tau+1$, just as in 1) for (63.). Repeating the process $\sigma-\tau$ times leads to the desired folding:
\[
 \pair{R^\lambda S^\sigma}{p_{\sigma+\lambda}}=\eps\pair{q_{n-\sigma-\lambda}S^\sigma}{T_{n-\tau}p_{\tau-\lambda}}.
\]
In the whole process only the normal-form property of $S$ was used. Viewing $T$ as a normal form gives the dual process, with complete reversal of the order, according to the formulas
\begin{align*}
\pair{q_{n-\sigma-\lambda}S^\sigma}{T_{n-\sigma}p_{\sigma-\lambda}}&=\eps\pair{T_{n-\sigma}R_\lambda}{p_{\sigma-\lambda}},\quad R_\lambda\sim q_{n-\sigma-\lambda}S^\sigma,\\
\pair{q_{n-\sigma}T^{\sigma-1}}{T_{n-\sigma}R_\lambda p_{\sigma-\lambda-1}}&=\eps\pair{T_{n-\sigma+1}R_\lambda}{p_{\sigma-\lambda-1}},\\
&\vdots\\
\pair{q_{n-\tau-1}T^\tau}{T_{n-\tau-1}R_\lambda p_{\tau-\lambda}}&=\eps\pair{q_{n-\sigma-\lambda}S^\sigma}{T_{n-\tau}p_{\tau-\lambda}}.
\end{align*}
3) Generation of cogredient foldings of defect $\lambda$ from cogredient foldings of defect $\lambda-1$ and fundamental foldings. The weight decrease for cogredient foldings of defect $\lambda$ is $2\lambda^2=2((\lambda-1)^2+2\lambda-1)$, giving the possibility of a decomposition into a cogredient folding of defect $\lambda-1$ and $2\lambda-1$ fundamental foldings. This is most easily found by first putting $n=2\lambda$. The only possible cogredient folding of defect $\lambda$ is $\pair{S^\lambda}{T_\lambda}=\pair{S^\lambda}{T_{n/2}}$; it contains one row $u$ and one row $x$ fewer than the folding of defect $\lambda-1$ of the same rows, $\pair{uS^\lambda}{T_\lambda x}$. Conversely, we compose the folding of defect $\lambda$ out of fundamental foldings which effect the disappearance of a row $u$ and a row $x$ (weight decrease $2(n-1)=2(2\lambda-1)$) and at the same time generate a new symbol row $R^\lambda$, together with a cogredient folding of defect $\lambda-1$. The simplest folding that does this is the following one of defect 1:
\[
 \pair{sT^\lambda}{\sigma p_\lambda}=\eps\pair{S^{2\lambda-1}}{R_{\lambda-1}p_\lambda}=\eps\pair{R^\lambda}{p_\lambda},
\]
where
\[
 R^{\lambda+1}=sT^\lambda,
 \quad s=S^1,
 \quad \sigma=S_1,
 \quad R^\lambda\sim\sigma R_{\lambda-1}.
\]
By (65.) and (64.) it is composed out of the $2\lambda-1$ fundamental foldings
\begin{equation}
 T^\lambda S^\lambda,S^{\lambda-1},\ldots,S^1;
 \qquad R^{\lambda+1}S^{\lambda+1},S^{\lambda+2},\ldots,S^{2\lambda-1}.
\tag{67}
\end{equation}
A folding of defect $\lambda-1$, taking account of (21.) and (62.), gives
\[
 \pair{uS^\lambda}{\sigma R_{\lambda-1}x}=\eps\pair{R^{\lambda+1}}{S_\lambda x}
 =\eps\pair{sT^\lambda}{S^\lambda x}=\eps\pair{S^\lambda}{T_\lambda}.
\]
The general case, for two rows $S^\rho,T^\rho$ and arbitrary $n$, is obtained directly from the special one. In place of (67.) one has the following $2\lambda-1$ fundamental foldings:
\begin{equation}
 T^\rho S^\rho,S^{\rho-1},\ldots,S^{\rho-\lambda+1};
 \qquad R^{\rho+1}S^{\rho+1},S^{\rho+2},\ldots,S^{\rho+\lambda-1}.
\tag{68}
\end{equation}
The first half of the foldings gives the form
\[
 \pair{q_{n-\rho-1}T^\rho}{S_{n-\rho+\lambda-1}p_{\rho-\lambda}},
\]
from which, by the substitutions
\[
 S_{n-\rho+\lambda-1}p_{\rho-\lambda}\sim w,
 \qquad T^\rho w=R^{\rho+1},
\]
and application of the second half of the foldings, one obtains
\[
 \pair{q_{n-\rho-\lambda}S^{\rho+\lambda-1}}{R_{n-\rho-1}p_\rho}.
\]
Substituting
\[
 q_{n-\rho-\lambda}S^{\rho+\lambda-1}\sim y,
\]
one obtains, by a cogredient folding of defect $\lambda-1$,
\begin{equation}
 \pair{q_{n-\rho-\lambda+1}S^\rho}{yR_{n-\rho-1}p_{\rho-\lambda+1}}.
\tag{69}
\end{equation}
We show that this is the desired folding of defect $\lambda$. Substitute
\[
 R_{n-\rho-1}p_{\rho-\lambda+1}=R_{n-\lambda}
\]
and observe that resolving $y$ gives
\[
 \pair{q_{n-\rho-\lambda+1}R^\lambda}{S_{n-\rho}y}
 =\eps\pair{q_{n-\rho-\lambda}S^{\rho+\lambda-1}}{S_{n-\rho}p'_{\rho-1}}=0.
\]
Thus by (20.) the value of (69.) is
\[
 \eps\pair{S^\rho R^\lambda}{p_{\rho+\lambda-1}y}
 =\eps\pair{q_{n-\rho-\lambda}S^{\rho+\lambda-1}}{[S^\rho R^\lambda]_{n-\rho-\lambda}p_{\rho+\lambda-1}}.
\]
This folding is of type (43.), hence equivalent to
\[
 \pair{q_{n-\rho-\lambda}S^\rho}{R_{n-\rho-1}p_{\rho-\lambda+1}}
 =\eps\pair{wT^\rho}{R_\lambda p_{\rho-\lambda+1}},
 \qquad R_\lambda\sim q_{n-\rho-\lambda}S^\rho.
\]
Here $R_\lambda$ is already of the desired final form; the expression corresponds to the form $\pair{sT^\lambda}{S_\lambda x}$ in the special case. Now, observing that resolving $w$ and $R_\lambda$ gives
\[
 \pair{wR^{n-\lambda}}{T_{n-\rho}p_{\rho-\lambda+1}}
 =\eps\pair{q'_{n-1}R^{n-\lambda}}{S_{n-\rho+\lambda-1}p_{\rho-\lambda}}
 =\eps\pair{q''_{n-\sigma-1}S^\rho}{S_{n-\rho+\lambda-1}p_{\rho-\lambda}}=0,
\]
we get by (21.) the value
\[
 \pair{wq_{n-\rho+\lambda-1}}{T_{n-\rho}R_\lambda}
 =\eps\pair{q_{n-\rho+\lambda-1}[T_{n-\rho}R_\lambda]^{n-\lambda}}{S_{n-\rho+\lambda-1}p_{\rho-\lambda}}
 \quad\text{equivalent to}\quad
 \pair{q_{n-\rho-\lambda}S^\rho}{T_{n-\rho}p_{\rho-\lambda}}.
\]
In the whole process only the normal-form property of $S$, not that of $T$, was used. Thus the cogredient folding of defect $\lambda-1$ used to generate (69.) can be replaced by $2\lambda-3$ fundamental foldings and a cogredient folding of defect $\lambda-2$, which in turn admits a decomposition into fundamental foldings, and so on. Analogously, the generation is obtained by using the normal-form property of $T$ and applying the dualistic process, that is, the fundamental foldings
\[
 S^\rho T^\rho,\quad T^{\rho-1},\ldots,T^{\rho-\lambda+1};
 \qquad R^{\rho-1}T^{\rho-1},\quad T^{\rho-2},\ldots,T^{\rho-\lambda+1},
\]
and a folding dual to (69.). The folding
\[
\pair{q_{n-\rho-\lambda}T^\rho}{S_{n-\rho}p_{\rho-\lambda}},
\]
equivalent to the preceding one, is then obtained.

Adding the decomposition obtained under 2), we have indeed obtained a generation of the most general foldings from fundamental foldings. It remains to show that this generation is unique, apart from permutation of the order, i.e. that to each folding there corresponds one and only one number $\alpha_\rho$ of fundamental foldings of type $\rho$, where $\rho=1,2,\ldots,n-1$.

Let $m_\rho$ be the degrees in the variable rows $p_\rho$ of the initial form $f$, and $l_\rho$ the degrees of the form produced by folding. Each fundamental folding of type $\rho$ transforms the degrees $m_{\rho-1},m_\rho,m_{\rho+1}$ respectively into $m_{\rho-1}+1,m_\rho-2,m_{\rho+1}+1$. Thus the $\alpha$ satisfy the system of equations
\begin{equation}
 m_\rho-l_\rho=-\alpha_{\rho-1}+2\alpha_\rho-\alpha_{\rho+1},
 \qquad (\rho=1,2,\ldots,n-1),
\tag{70}
\end{equation}
where $\alpha_0$ and $\alpha_n=0$ are to be put, and whose determinant has value $n$.\footnote{If $\Delta_n$ denotes the $n$-rowed determinant, then the recurrence formula holds: $\Delta_n=2\Delta_{n-1}-\Delta_{n-2}$; that is, $\Delta_n-\Delta_{n-1}=\Delta_{n-1}-\Delta_{n-2}=\cdots=\Delta_2-\Delta_1=1$, since $\Delta_2=3$, $\Delta_1=2$.} The $\alpha$ are thereby uniquely determined, and indeed, by Theorem VIII, as positive integers; this gives conditions for the differences $m_\rho-l_\rho$.

The results of this paragraph hold by virtue of relations (62.); however, if the relations (62.) do not hold, the final expressions are by no means equivalent to the initial expressions. The definition of equivalence given in \S\ 6 also admits the formulation:\footnote{Cf. the consideration at the end of the next paragraph.} ``Two forms are equivalent if and only if they differ by forms of higher folding, i.e. by forms that exhibit a greater weight decrease.'' This greater weight decrease always occurs when the two rows $q_{n-\sigma-i},p_{\tau-\lambda}$ entering the folding are really variable rows, as happens in (41.), (42.), in the foldings appearing in this paragraph as equivalent, and also in the equivalent forms of Theorem VII. But it need not occur when these rows are derived from symbol rows by substitution, as was the case for the forms of this paragraph which could be expressed in terms of one another by virtue of (62.). One easily sees that the vanishing forms even have, in part, a higher total weight.

\section*{\S\ 9. Form rows. Reduction theorems.}
We now define the form row introduced in \S\ 6 (end) somewhat more sharply as ``the totality of invariant formations of a given initial form that are linear in the coefficients, i.e. the totality of forms that arise from the initial form by folding in itself, arranged according to higher forms.''

To explain the higher forms, suppose by Theorem VII that the initial form is given as a product of two normal forms $S\cdot T$. By a higher form we mean forms with higher folding in itself, and among forms with the same folding in itself, specially normalized ones. More precisely: suppose the form $A$ has arisen by $\alpha_\rho$ fundamental foldings of type $\rho$, and the form $B$ by $\beta_\rho$ such foldings. Then $B$ is higher than $A$:
\begin{enumerate}
\item if, in the system of inequalities $\beta_\rho\ge\alpha_\rho$, $\rho=1,2,\ldots,n-1$, the inequality sign holds for at least one value of $\rho$;\footnote{If in the system of inequalities the sign $>$ holds in part and the sign $<$ in part, then the forms are not comparable, since a form arising from another by folding always has a positive value system of fundamental foldings, hence in our case positive or vanishing values $\beta_\rho-\alpha_\rho$.}
\item if the equality sign holds for all values $\rho$, but fewer symbol rows have been combined into a single matrix product. This normalization proves necessary because of the reductions in \S\ 8. Accordingly, for example, $\pair{S^{n-1}}{T_{n-1}}$ is higher than $\pair{S^{n-1}}{[S^1T^\rho]_{n-\rho-1}p_\rho}$;
\item if the equality sign holds for all values $\rho$ and the number of symbol rows combined into one matrix product is the same, then $B$ is made higher than $A$ by a normalization arbitrary in itself but fixed once and for all. This normalization is always possible because of the finite number of forms belonging to a value system $\alpha$; it can be achieved, for instance, by privileging the one form $S$ (larger number of symbol rows $S$, larger sum of the upper weight indices of the rows $S$, and so on).
\end{enumerate}
It is useful to think of the form row as arranged as an $(n-1)$-dimensional manifold of lattice points, where the $n-1$ directions correspond to the $n-1$ fundamental foldings. To every form there then corresponds one and only one value system $\alpha$, i.e. a positive integral lattice point, while the different forms belonging to a value system $\alpha$ are uniquely determined in their order by the above normalization. An arbitrary form of the form row gives the initial form of a new form row, contained as a part of the total form row; indeed it is contained as the part consisting of the forms obtained from the new initial form by proceeding in the $n-1$ positive directions.

By virtue of Theorem VIII and the general theory of series expansions, exactly the reduction theorems of the binary and ternary domains (cf. dissertation, \S\ 3) hold for form rows. We shall briefly derive the principal theorem. We first state the definition:

``In a given system of forms, suppose a certain finite number of forms has been combined into a module. We call a form reducible if it can be expressed by forms that have invariants as factors, or by `higher forms'; that is, by higher forms of the total form row to which the form belongs, or by forms that contain the symbols of the module in higher order. A reducible form row is called a reducer.'' Then:

\emph{Theorem IX:} If the initial form of a form row is reducible by having arisen through folding with a reducer, then the entire form row arising from this initial form is reducible.

For the proof, observe that a form $h$ arising by folding a form $f$ with a second form $g$ can, because of (45.), be regarded as a form $f$ containing several cogredient variable rows $p_\rho,p'_\rho$. Such forms can, by the general theory of series expansion, be represented as polars of the elementary covariants of $f$, by applying the series expansion successively to each pair of cogredient variable rows $p_\rho,p'_\rho$. By (38.), the elementary covariants that occur are the forms generated by cogredient folding. Since, however, the order in which the series expansion is applied is still arbitrary, we may in particular choose an order corresponding to the generation of the general foldings from fundamental foldings. The resulting system of elementary covariants is then identical with the form row $f$; the forms that arise by cogredient folding of higher defect are arranged among those that arise by fundamental foldings. But one also obtains polars of the form row $f$ for any order of the series expansion, to which a second system of elementary covariants may correspond. Since the form row exhausts the totality of invariant formations, every form of the second system can be expressed linearly by polars of the form row, and indeed by polars of those forms that show the same or a greater weight decrease. The latter follows because (cf. \S\ 7) the polars can be regarded as simultaneous invariants of a form row $H$ (52.), with the degree numbers corresponding to the form $h$, and of the form row $f$. Every folding of $H$ in itself corresponds to a weight decrease which, after substitution (11.) is performed, is transferred to the symbol rows and hence to the forms of $f$.\footnote{This consideration also underlies the second formulation of the equivalence concept (\S\ 8, end). Further explanations of the different systems of elementary covariants in the ternary domain are found in Study, loc. cit., II, \S\ 7, Theorems 7) and 8). By our Theorem VII the same results also hold for $n$ variables.}

An arbitrary form of the form row $h$ has arisen by folding $h$ in itself; that is, on the one hand by a higher folding of $g$ with $f$, and on the other by folding $f$ or $g$ in itself. In both cases the series expansion leads to polars of the form row $f$, just as was shown above for the initial form $h$. If in particular $f$ is a reducer, then an expansion in polars of the reducer results; hence the form row $h$ becomes reducible.

Applications of the theorem to the reduction of complete systems of forms follow analogously to the binary and ternary domains.

\end{document}
